\newpage
\section{M. Carlsen - H. Nakamura, Wijk aan Zee 2011}
\begin{multicols}{2}
    \newchessgame
    \chessgameinfo{Wijk aan Zee}{M. Carlsen}{H. Nakamura}{}{2011}{1-0}

    \mainline[level=1]{
        1. e4 c5 2. Nf3 d6 3. d4 cxd4 4. Nxd4 Nf6 5. Nc3 a6 6. Be2 e5 7. Nb3 Be7 8. Be3 O-O
    }

    \chessboard

    \mythoughts{
        Normally White castles queenside and launches a pawn storm on that flank.
        \variation{9. Qd2} is natural, but \variation{9... Ng4}, exchanging the bishop,
        is a nuisance. We want to drive the f6 knight away with g4 and g5, so there is no
        reason to delay.
    }

    \myreflection{
        The idea above is good. By the way, the computer does not like Black's last move,
        castling, perhaps because it is too committal.
    }

    \mainline[level=1]{
        9. g4 Be6 10. g5 Nfd7
    }

    \mythoughts{
        Several natural moves are available: \symqueen d2, O-O-O and h4.
        There is no reason to delay the queen move, since I will certainly
        castle queenside.
    }

    \myreflection{
        There is an important difference between \variation[invar]{11. h4} and \variation{11. Qd2}.
        After \variation{11. Qd2 a5 12. f4 exf4 13. Bxf4 Nc6 14. Nd4? Nxd4 15. Qxd4 Bxg5} Black wins a pawn,
        whereas after \variation{11. h4 a5 12. f4 exf4 13. Bxf4 Nc6 14. Nd4? Nxd4 15. Qxd4} the g5 pawn is still protected.
    }

    \mainline[level=1]{
        11. h4 Nb6 12. Qd2 N8d7
    }

    \mythoughts{
        I could launch a brute-force attack: h5, g6, then double my rooks on
        the g-file. My dark-squared bishop would go to h6, and the queen could
        also swing to the g-file. That would give four attackers on the g7 pawn.

        Before committing to these moves, let me first play \variation{13. O-O-O}.
    }

    \myreflection{
        \variation[invar]{
            13. O-O-O
        } is in fact the move recommended by Stockfish. I realize that the plan with f4
        is new to me.
    }

    \mainline[level=1]{
        13. f4 exf4
    }

    \mythoughts{
        I see no reason to play \variation{14. Bf2} or \variation{14. Bg1}.
        The d4 square is a perfect outpost for my b3 knight, so I decide to play \variation{14. Bxf4}.
    }

    \mainline[level=1]{
        14. Bxf4 Ne5
    }

    \mythoughts{
        I still need to castle, so I play \variation{15. O-O-O}.
    }

    \myreflection{
        Of course castling is a good move. I like positions where the natural
        moves are also the good moves.
    }

    \mainline[level=1]{
        15. O-O-O Rc8
    }

    \mythoughts{
        Heading into an endgame with \variation{16. Bxe5 dxe5 17. Qxd8 Rcxd8 18. Rxd8 Rxd8 19. Rd1 Rxd1 20. Kxd1}
        is inconsistent with my original kingside attack. The endgame is bad for White because Black
        has the bishop pair and the h4-g5 pawns are weak.

        The exchange \variation{16. Nd5 Nxd5 17. exd5 Bf5 18. Bd3 Bxd3 19. Qxd3} is also unsatisfactory.
        Black's weakness on d6 has disappeared, and there is no static advantage left to play for.

        I notice that my c2 pawn is weak. Soon Black will play \symrook xc3, \symknight a4, \symqueen c7,
        \symrook c8, d5, \symbishop f5. I may have enough defenders, but defending postpones
        my own attack. If my king stays on c1, playing c3 to defend would cut the king off and prevent it
        from escaping to the safer corner on b1 or a1.

        I want to play \variation{16. Kb1}.
    }

    \myreflection{
        The calculation \variation[invar]{16. Bxe5 dxe5 17. Qxd8 Rcxd8 18. Rxd8 Rxd8 19. Rd1 Rxd1 20. Kxd1}
        is completely correct.

        In \variation[invar]{
            16. Nd5 Nxd5 17. exd5 Bf5 18. Bd3?? Nxd3
        } I simply lose a bishop!

        Magnus plays so many \symking b1 and \symking a1 moves. Here is another example, Carlsen vs Li Chao, Qatar Masters 2015:
    }


    \chessboard[
        setfen=r1b1qrk1/ppp1p1bp/1n4p1/4Pp2/1n1P4/2N1BP1N/PP1Q2PP/2KR1B1R w - - 3 12
    ]

    \myreflection{
        He also played \symking b1 and then later \symking a1. Note that in both the current
        game and that one, Black's light-squared bishop is on e6, so moving the king to a1
        is a useful prophylactic measure.
    }


    \mainline[level=1]{
        16. Kb1 Qc7
    }

    \mythoughts{
        I remember a game between Carlsen and Li Chao at the Qatar Masters. Li also had
        a bishop on e6 in a Gr\"unfeld Defence, and Carlsen played \symking a1 first to avoid any potential
        checks from that bishop. I decide to play \variation{17. Ka1}.
    }

    \myreflection{
        According to the computer, \variation[invar]{
            17. Ka1
        } is just as good as the text move. I like it
        because mixing kingside action with queenside prophylaxis requires
        more thought. Why not finish all the prophylaxis first?
    }

    \mainline[level=1]{
        17. h5 Rfe8
    }

    \mythoughts{
        Of course I would love to play g6, hxg6 and then swing my queen to the kingside.
        Before that, I must move my e2 bishop somewhere. Finding a good square is
        difficult and costs a tempo. What if I play \variation{18. Ka1}, a move that is always useful?

        Black must then play \variation{18... Nbc4}, and I can capture the knight with my bishop
        without wasting a tempo searching for a good square for it.
    }

    \myreflection{
        Saving a tempo on the bishop move is very interesting. \variation[invar]{
            18. Ka1 Nec4 19. Bxc4 Nxc4 20. Qd3
        } is necessary, preventing the \variation{20... Nxb2} idea.
    }

    \mainline[level=1]{
        18. Ka1 Bf8
    }

    \mythoughts{
        Black's last move is cunning: he refuses to help White save a tempo.

        Let me try some schematic thinking. White will play g6, since h6 makes no sense.
        White's rook and queen belong on the h-file, because Black already defends
        the g-file well. To prevent Black's king from sneaking away, one rook should
        stay on the f-file. Before swinging the queen to h2, the bishop must move first;
        d1, vacated by the rook, is a perfect square, since from there the bishop also
        protects the critical c2 square.

        I decide to play \variation{19. Rdf1}.
    }

    \myreflection{
        The schematic-thinking idea is good, but I still have to calculate concretely:
        \variation[invar]{
            19. Rdf1 Nbc4 20. Bxc4 Nxc4 21. Qd3 Na3 22. Rf2 b5 23. Rc1 Nc4 24. g6 fxg6 25. hxg6 h6
        }. Note Black's typical knight maneuver: c4-a3, targeting the c2 pawn.
        Black parries White's attack by generating enough counterplay on the queenside.

        In practice my idea is almost as good as the text move. The text move
        gives Black the chance to go wrong, while mine does not.

        The computer recommends \variation[invar]{
        19. Qe1 Nac4 20. Qg3
        }: the g3 square belongs to the queen, as we will see.
    }

    \mainline[level=1]{
        19. Nd4 Qc5
    }

    \variation[invar]{
        19... Nec4 20. Bxc4 Nxc4 21. Qd3 Qa5 22. Nb3 Qc7 23. Qg3 Na3 24. Rh2 Qc4 25. Rf2
    } is better. Note that Black places his knight on a3 only after White's queen moves to g3,
    forcing White to make a wasted rook move to h2.

    \chessboard

    \mythoughts{
        I was intrigued by the idea \variation{20. Nf5 Rcd8 21. Nh6+}.
        Black can simply play \variation{21... Kh8}, and I don't see
        how to proceed with my knight on h6, since the e6 bishop protects f7.

        Black can even play \variation{20... Qa5} and White will lose material.

        I rejected this idea because Black's e6 bishop is too strong. \variation{20. Nxe6} is natural:
        \variation{20... Rxe6 21. Bxe5 Qxe5 22. Bg4} and White wins material;
        on \variation{20... fxe6 21. g6 h6} I can break through along the f-file, using f7 as an
        outpost.
    }

    \myreflection{
        My calculation was wrong. After \variation[invar]{
            20. Nxe6 fxe6 21. g6 Nc4
        } White must exchange queens, since Black threatens to take the b2 pawn with the knight, and once the queens come off there is no attack left.

        I do not understand why the natural text move did not occur to me.
        In computer-science terms, a variation is a tree. If my calculation is wrong, either
        I did not go deep enough or I missed some moves at the start. I realize that missing
        moves at the start is by far the more likely failure, so I must
        keep asking myself whether I have any further options.
    }

    \mainline[level=1]{
        20. g6 Nec4 21. Bxc4 Nxc4
    }

    \mythoughts{
        I thought \variation{22. Qh2?} was the only move.
        However, after \variation{22... Nxb2 23. Kxb2 Qb4+ 24. Kc1 Qxc3}
        my king is in danger and my attack has not yet materialized.
        The c3 knight must be protected, and there is only one way to do it.
    }

    \myreflection{
        The use of process of elimination here is very good. I must apply it more often to avoid
        blunders!

        It is also instructive to see that Black does not take the pawn immediately, but inserts an
        intermediate move first. I wish I would never miss intermediate moves in my own calculations.
    }

    \mainline[level=1]{
        22. Qd3 fxg6
    }

    \mythoughts{
        Black threatens \variation{23... d5 24. exd5 Nxb2 25. Kxb2 Qa3+ 26. Kb1 Rxc3}, which is clearly
        winning.

        I am reluctant to exchange the e6 bishop for my d4 knight. Black's king often
        escapes via f7, and with my knight on g5 that route is blocked.

        On the other hand, Black also intends to mate in a few moves with the queen on a5, sacrificing
        on c3 with his rook and mating on a2, as mentioned above. The e6 bishop plays a vital role,
        while I would need more tempi to mate with my knight on g5.

        Black also intends \symqueen b4, threatening mate on b2; this can be
        parried with \symbishop c1. Note that \symbishop c1 is not merely defensive, since it opens
        the f-file, cutting off Black's king.

        I therefore want to play \variation{23. Nxe6}. One possible line is
        \variation{23... Qb4 24. Bc1 Rxe6 25. Qh3 Rce8 26. hxg6 h6 27. Nd5}: White has more attackers
        than Black, and should be better.
    }

    \myreflection{
        I realize I am calculating concrete moves far too late, only after Black has
        very concrete ideas of his own. I should have thought more thoroughly on move 19, four moves ago.

        My calculation above was fine: White is indeed much better after
        \variation[invar]{23... Qb4 24. Bc1 Rxe6 25. Qh3 Rce8 26. hxg6 h6 27. Nd5}.
        An even simpler way is
        \variation[invar]{23... Qb4 24. b3}, winning material.

        Black must capture the knight directly, without any intermediate move:
        \variation[invar]{
            23... Rxe6 24. Qh3 Rce8 25. hxg6 h6 26. Nd5! Rxg6 27. Qf5! Ne5 28. Bxe5 Rg5 29. Qf3 Rgxe5 30. Rdf1
        }. I am not sure I could calculate this line on my own.

        The computer does recommend my move over the text move.
    }

    \mainline[level=1]{
        23. hxg6 h6
    }

    \mythoughts{
        The idea is quite simple: I sacrifice on h6 and mate with the queen.
        Black has at most two defenders for the h6 pawn, and White can bring more,
        so White should be winning provided Black cannot checkmate first.

        White wants to play \symqueen h3, since from the third rank the queen also keeps an eye
        on the important c3 knight. This is currently not possible because of the e6 bishop,
        so removing it is a natural idea.

        \variation{24. Nxe6 Rxe6}: Black wins a tempo by attacking the g6 pawn, and can then play
        ...d5 and \symrook b6 with a strong attack.

        \variation{24. Nd5 Bxd5 25. exd5 Qxd5 26. Rh4 Be7 27. Qh3}: White will sacrifice
        on h6 and mate.
        \variation{25... Qb4 26. Nb3}: White is safe.

        Black does not have to take the d5 knight, but White will simply capture the e6 bishop
        himself. With the knight on d5, the ...d5 break is no longer available for Black, as noted before.
    }

    \myreflection{
        I should have calculated more concretely. After \variation[invar]{24. Nxe6 Rxe6 25. Rdg1 Qb6 26. Bc1 d5 27. exd5}
        Black has no attack and White is clearly better.

        \variation[invar]{
            24. Nd5 Bxd5 25. exd5 Qxd5 26. Rh4 Be7 27. Qh3? Bxh4 28. Qxh4
        }: White can still sacrifice on h6, but there will be no mate on h7, since Black
        has enough defenders.

        Comparing \variation{24. Nxe6} and \variation{24. Nd5}, White must decide
        which knight to keep. The c3 knight is the better choice, since it can go to d5. Such
        a good knight against the bad f8 bishop is useful not only in the middlegame
        but also in the endgame.

        I missed the text move completely. The queen can also go to g3. After a sacrifice
        on h6, White can then push the g6 pawn and promote.
    }

    \mainline[level=1]{
        24. Qg3 Qb6
    }

    \mythoughts{
        Black threatens checkmate on b2. \variation{25. Rb1} is out of
        the question because it is too passive.
        \variation{25. Nb3} is not good either: Black can always drive
        the knight away with a5 and a4, and besides, the knight is better
        placed on d4 for the attack.
        Therefore, I want to play \variation{25. Bc1}, defending b2
        and at the same time opening the f-file for the attack.
    }

    \myreflection{
        A very effective thinking process.
    }

    \mainline[level=1]{
        25. Bc1 Qa5
    }

    \mythoughts{
        Black threatens \variation{26... Nxb2}, \variation{27... Rxc3} and \variation{28... Qxa2#}.

        White has his own threat: \variation{26. Rdf1}, \variation{27. Rxf8}, \variation{28. Bxh6},
        \variation{29. Bxg7}, \variation{30. Qh3} and \variation{31. Qh7#}.

        Black is clearly faster, because his moves are more forcing and because he needs fewer of them,
        so it makes sense for White to defend first.

        \variation{26. Nd5} loses to \variation{26... Nb6}.
        \variation{26. b3 Na3 27. Bxa3 Qxa3}: Black is better, because White loses an important attacker
        who could otherwise sacrifice on h6.
        I do not want to play \variation{26. Nxe6}, exchanging the pieces, if I have better options,
        since Black's c8 rook can join the attack via c6 and b6.

        \variation{26. Rdf1 Nxb2? 27. Rxf8+ Kxf8 28. Nxe6+ Rxe6 29. Bxb2}: White is better.
        Compared with \variation{26. Nxe6}, White gains significantly more by removing
        the f8 bishop from the board.

        I want to play \variation{26. Rdf1}.
    }

    \myreflection{
        I made a big mistake here: ``\variation{26. Nd5} loses to \variation{26... Nb6}.''
        In fact \variation[invar]{26. Nd5 Nb6?? 27. Nxe6 Rxe6 28. Qg4 Rce8 29. Qf5 Rxg6} and White is crushing.

        ``I do not want to play \variation{26. Nxe6}, exchanging the pieces, if I have better options,
        since Black's c8 rook can join the attack via c6 and b6.''
        After \variation[invar]{26. Nxe6 Rxe6 27. Qh3 Rce8 28. Nd5} the b6 square is no longer available
        for Black's rook. White is indeed much better: he has a safer king, and his knight is better than the bishop.

        I should be happy to have found the text move. Its intention, however, is not to sacrifice
        on f8, as we will see.
    }

    \mainline[level=1]{
        26. Rdf1 Ne5
    }

    \chessboard 

    \mythoughts{
        Black still has the same threat.
        After \variation{27. Nxe6 Rxe6} I would have to protect the g6 pawn with a rook,
        but both of my rooks are targeting the important f8 and h6 squares
        and cannot be spared.

        \variation{27. Nd5 Bxd5 28. exd5 Qxd5 29. Nf5}: White threatens
        \variation{30. Nxh6} and \variation{31. Rxf8}, and the g-pawn will promote soon.

        I decide to play \variation{27. Nd5}.
    }

    \mainline[level=1]{
        27. Nd5 Bxd5 28. exd5 Qxd5 29. Bxh6 gxh6
    }

    \myreflection{
        I didn't realize that I could sacrifice on h6.
    }

    \chessboard

    \mythoughts{
        I still believe I must move my knight, since it is hanging.
    }

    \myreflection{
        \variation[invar]{
            30. Nf5
        } is also not a bad move, because White is still winning. The sacrifice on f8 is, however,
        not correct. White threatens to push the g-pawn to g7, and his advantage is so large that he
        can afford to miscalculate here.
    }
    \mainline[level=1]{
        30. g7 Be7
    }

    \mythoughts{
        The next move is obvious: \variation{31. Rxh6}, threatening mate on h8.
    }

    \mainline[level=1]{
        31. Rxh6 Nf7
    }

    \chessboard

    \mythoughts{
        There is a forced mate:
        \variation{32. Rxf7 Kxf7 33. Qf4+ Kg8 34. Rh8+ Kxg7 35. Qh6+ Kf7 36. Rh7+ Kg8 37. Qg7#}.
    }

    \myreflection{
        \variation{32. Rxf7? Qg5! 33. Rf8+ Rxf8 34. gxf8=Q+ Rxf8 35. Qxg5+ Bxg5} throws away the win!

        I played a good game, built up a winning advantage, and then missed the win. Missing
        an intermediate move like this is painful.
    }
    \mainline[level=1]{
        32. Qg6 Nxh6 33. Qxh6 Bf6 34. Qh8 Kf7 35. g8=Q Rxg8 36. Qxf6 Ke8 37. Re1
    }
\end{multicols}

\subsection*{My Reflection}
I am happy that I found good moves through logical reasoning, without deep calculation:
I drew on similar games from memory to extract general ideas, used schematic thinking
to land on the same move as Carlsen, and applied process of elimination to rule out bad moves.

I also realize that I made mistakes in my calculations. The remedy is to keep asking myself:
did I miss any move at the start? Even when my calculation is not deep enough, this question
helps me eliminate the worst options.

Another lesson is that my assessments must be backed up by concrete calculation. This does not
mean I have to calculate everything; but I must at least support my assessment with a few
sample lines.
